

\medskip
\phantomsection\label{prop:sga3-VIB-10.12}%
\noindent\textbf{Proposition 10.12.---}
\begingroup\itshape
Let \(S\) be an integral scheme, \(G\) an \(S\)-group of finite
presentation, and \(A\) and \(B\) two group subschemes of \(G\) of
finite presentation, with smooth generic fiber, such that \(A\) has
connected generic fiber \textup{(}respectively, such that \(A\) and
\(B\) are invariant in \(G\)\textup{)}. Then there exist a nonempty
open subset \(S'\) of \(S\) and a closed group subscheme \(D'\) of
\[
  G'=G|_{S'},
\]
of finite presentation, with smooth fibers, such that \(D'\) has
connected fibers \textup{(}respectively, is invariant in
\(G'\)\textup{)}, and such that \(D'\) represents the sheaf
associated, for the topology \textup{(fppf)} or \textup{(fpqc)}, with
the presheaf of groups of commutators of \(A\) and \(B\) in \(G\). In
particular, for every \(s\in S'\),
\[
  D'_s=(A_s,B_s),
\]
with the notation of \hyperref[def:sga3-VIB-7.2-vii]{7.2~\textup{(vii)}}.
\par\endgroup

\emph{Proof.}
Let \(\eta\) be the generic point of \(S\), and put
\[
  H_\eta=(A_\eta,B_\eta).
\]
By \hyperref[prop:sga3-VIB-7.8]{7.8} \textup{(}respectively,
\hyperref[concl:sga3-VIB-7.3-v]{7.3~\textup{(v)}}\textup{)}, \(H_\eta\) is
connected \textup{(}respectively, invariant in \(G_\eta\)\textup{)}. By
\hyperref[par:sga3-VIB-10.1]{10.1} and
\hyperref[par:sga3-VIB-10.10-a]{10.10~a)}, there exist a nonempty open
subset \(S'\) of \(S\) and an \(S'\)-group subscheme \(D'\) of finite
presentation and closed in \(G'\), with geometrically reduced fibers
\textup{(}EGA~IV\(_3\), 9.7.7 and 9.3.3\textup{)}, having \(H_\eta\) as its
fiber at the point \(\eta\), and such that \(D'\) has connected fibers
\textup{(}that is, geometrically connected fibers
\textup{(\hyperref[prop:sga3-VIA-2.1.1]{VI\(_A\), 2.1.1}))}
\textup{(}respectively, is invariant in \(G'\)\textup{)}
\textup{(}EGA~IV\(_3\), 9.7.7 and 9.3.3 \textup{(}respectively,
\hyperref[lem:sga3-VIB-10.4]{10.4}\textup{))}.\footnote{Translator's note: the Gille--Polo re-edition
(2011), Expos\'e~VI\textsubscript{B}, also cites \hyperref[lem:sga3-VIB-10.3]{10.3} here, which gives the group
structure of \(D'\).} Moreover, we saw in the
course of the proof of \hyperref[prop:sga3-VIB-7.8]{7.8} that there is an
integer \(n\) such that
\phantomsection\label{form:sga3-VIB-10.12-generic-image}%
\[
  \nu_\eta^n
  \left((A_\eta\times_{\kappa(\eta)}B_\eta)^n\right)
  =H_\eta,
\]
where \(\nu_\eta:A_\eta\times_{\kappa(\eta)}B_\eta\to G_\eta\) is defined as in
\hyperref[def:sga3-VIB-7.2-vii]{7.2~\textup{(vii)}}. By the same formulas as in
\hyperref[def:sga3-VIB-7.2-vii]{7.2~\textup{(vii)}} and
\hyperref[concl:sga3-VIB-7.1-ii]{7.1~\textup{(ii)}}, we can define the morphisms
\phantomsection\label{form:sga3-VIB-10.12-nu-prime}%
\[
  \nu':A'\times_{S'}B'\longrightarrow G',
  \qquad
  {\nu'}^n:(A'\times_{S'}B')^n\longrightarrow G',
\]
where \(A'=A|_{S'}\) and \(B'=B|_{S'}\); and indeed
\[
  {\nu'}^n\otimes_{S'}\kappa(\eta)=\nu_\eta^n.
\]
Consequently \textup{(}EGA~IV\(_3\), 8.8.3, 8.10.5 and 11.2.6\textup{)}, we
can choose \(S'\) so that the morphism \({\nu'}^n\) is flat and factors
through \(D'\), and so that the morphism \((A'\times_{S'}B')^n\to D'\)
deduced from \({\nu'}^n\) is surjective.\footnote{Translator's note: here, as in the Gille--Polo
re-edition (2011), the text asserts that \({\nu'}^n\) can be made flat. This is not justified for an integer \(n\) as
in 7.8, and can fail: for \(A=B=G=\mathrm{SL}_2\) over an algebraically closed field, every element is a
commutator, so \(n=1\) will do, but \(\nu\) is not flat at \((e,e)\), its fiber there (the pairs of commuting
elements) having dimension \(4\), not \(6-3=3\) (EGA~IV\(_2\), 6.1.2). Lemma~7.5 only needs
\({\nu'}^n\colon(A'\times_{S'}B')^n\to D'\) to be flat at some point above each maximal point of each fiber of
\(D'\), and this can be arranged: \(\nu_\eta^n\) is flat above a dense open subset of \(H_\eta\)
(EGA~IV\(_2\), 6.9.1), so, after shrinking \(S'\) (\hyperref[par:sga3-VIB-10.1]{10.1} and EGA~IV\(_3\), 9.5.5),
\({\nu'}^n\) is flat above an open subset of \(D'\) whose complement has fibers of smaller dimension than those of
\(D'\). The conclusion of 10.12 stands.} Then
\textup{(\hyperref[lem:sga3-VIB-7.5]{7.5})} the morphism
\phantomsection\label{form:sga3-VIB-10.12-covering-map}%
\[
  {\nu'}^{2n}
  =\mu\circ\bigl({\nu'}^n\times_{S'}{\nu'}^n\bigr):
  (A'\times_{S'}B')^{2n}\longrightarrow D'
\]
is a covering for the fppf topology \textup{(}and a fortiori for the
fpqc topology\textup{)}.\footnote{Translator's note: the 1970 text has \((A'\times_{S'}B')^{n+1}\) and
\({\nu'}^{n+1}\). Lemma~7.5, applied to \(f={\nu'}^n\), which is flat and surjective onto \(D'\), gives the covering
\(\mu\circ(f\times_{S'}f)\), defined on \((A'\times_{S'}B')^{2n}\); the Gille--Polo re-edition (2011),
Expos\'e~VI\textsubscript{B}, corrects accordingly.} Therefore \(D'\) represents the sheaf associated,
for the fppf or fpqc topology, with the presheaf of commutators of \(A\)
and \(B\) in \(G\).\footnote{Translator's note: this is 7.6 in the relative form given in the Addition
after \hyperref[prop:sga3-VIB-7.6]{7.6} (the equivalence of its conditions (i) and (ii)), which the Gille--Polo
re-edition (2011), Expos\'e~VI\textsubscript{B}, cites here, applied to \(f=\nu'\colon A'\times_{S'}B'\to D'\):
\(\nu'\) factors through \(D'\) because \({\nu'}^n\) does, \(\nu'(x)\) being \({\nu'}^n(x,e,\ldots,e)\); and the
morphism \(f^\infty\) of that Addition is a covering, since \(f^{2n}={\nu'}^{2n}\) is one.} Moreover, for every
\(s\in S'\), the morphism
\[
  (A_s\times_{\kappa(s)}B_s)^n\longrightarrow D'_s
\]
deduced from \(\nu_s^n\) is surjective; hence
\textup{(\hyperref[prop:sga3-VIB-7.6]{7.6})}, \((A_s,B_s)\) represents the
sheaf of commutators of \(A_s\) and \(B_s\) in \(G_s\) \textup{(}for the
fppf or fpqc topology\textup{)}, whence \(D'_s=(A_s,B_s)\).\footnote{Translator's note: the
Gille--Polo re-edition (2011), Expos\'e~VI\textsubscript{B}, details this last step. In short: the closed subgroup
\(D'_s\) of \(G_s\) contains the image of \(\nu_s\), since \(D'\) contains that of \(\nu'\), hence contains
\((A_s,B_s)\); conversely \(\nu_s^n\) maps onto \(D'_s\) and has its image in \((A_s,B_s)\); as \(D'_s\) is reduced,
being smooth, \(D'_s=(A_s,B_s)\).}

\par\medskip\noindent
\emph{Addition} (after the Gille--Polo re-edition (2011), Expos\'e~VI\textsubscript{B}; our formulation).
The re-edition adds here a corollary (its 10.12.1), which it uses in its proof of
\hyperref[prop:sga3-VIB-8.4]{8.4}. Let \(S\) be an integral scheme and \(G\) an \(S\)-group of finite
presentation with smooth fibers. For every point \(t\) of \(S\), put \(K^0_t=G_t\) and, for \(p\geq1\),
\(K^p_t=(K^{p-1}_t,K^{p-1}_t)\) \textup{(}respectively \(K^p_t=(G_t,K^{p-1}_t)\)\textup{)}, with the notation of
\hyperref[def:sga3-VIB-7.2-vii]{7.2~(vii)}; these are smooth closed subgroups of \(G_t\). Then, for every
\(n\geq0\), there are a nonempty open subset \(S_n\) of \(S\) and a group subscheme \(D_n\) of \(G|_{S_n}\),
invariant, of finite presentation and with smooth fibers, such that \((D_n)_s=K^n_s\) for every \(s\in S_n\).
Indeed, take \(S_0=S\) and \(D_0=G\). If \(S_n\) and \(D_n\) have been found, apply
\hyperref[prop:sga3-VIB-10.12]{10.12} over the integral scheme \(S_n\), in the case where both group subschemes
are invariant, to \(A=B=D_n\) \textup{(}respectively \(A=G|_{S_n}\), \(B=D_n\)\textup{)}: it gives a nonempty
open subset \(S_{n+1}\) of \(S_n\) and an invariant group subscheme \(D_{n+1}\) of \(G|_{S_{n+1}}\), of finite
presentation and with smooth fibers, such that \((D_{n+1})_s=((D_n)_s,(D_n)_s)=K^{n+1}_s\)
\textup{(}respectively \((D_{n+1})_s=(G_s,(D_n)_s)=K^{n+1}_s\)\textup{)} for every \(s\in S_{n+1}\).

\medskip
\phantomsection\label{cor:sga3-VIB-10.13}%
\noindent\textbf{Corollary 10.13.---}
\begingroup\itshape
Let \(S\) be an integral scheme with generic point \(\eta\), \(G\) an
\(S\)-group of finite presentation, and \(H\) an \(S\)-group subscheme of
\(G\), of finite presentation, with smooth fibers and invariant in \(G\).
Suppose that \((H_\eta,H_\eta)=H_\eta\) \textup{(}respectively, that
\(G_\eta\) is smooth and that \((G_\eta,H_\eta)=H_\eta\)\textup{)}.\footnote{Translator's
note: the 1970 text does not assume \(G_\eta\) smooth. The commutator group \((G_\eta,H_\eta)\) of 7.2~(vii) is only
defined when \(G_\eta\) is smooth, and 10.12, applied with \(A=G\), requires it; \(G\) is then smooth over a nonempty
open subset of \(S\) (\hyperref[par:sga3-VIB-10.1]{10.1} and \hyperref[par:sga3-VIB-10.10-a]{10.10~a)}), so that
\((G_s,H_s)\) is defined there. The Gille--Polo re-edition (2011), Expos\'e~VI\textsubscript{B}, adds this hypothesis,
and assumes only that the generic fiber of \(H\) is smooth, which suffices for the same reason.} Then there exists a
nonempty open subset \(S'\) of \(S\) such that, for every \(s\in S'\), \((H_s,H_s)=H_s\) \textup{(}respectively,
\((G_s,H_s)=H_s\)\textup{)}.
\par\endgroup

This follows immediately from \hyperref[prop:sga3-VIB-10.12]{10.12} and
EGA~IV\(_3\), 8.8.2.5.\footnote{Translator's note: the Gille--Polo re-edition (2011),
Expos\'e~VI\textsubscript{B}, details the argument. In short: 10.12, applied in its invariant case to \(A=B=H\)
\textup{(}respectively \(A=G\), \(B=H\)\textup{)}, gives a nonempty open subset \(S'\) of \(S\) and a closed group
subscheme \(D'\) of \(G|_{S'}\), of finite presentation, with \(D'_\eta=(H_\eta,H_\eta)=H_\eta\)
\textup{(}respectively \(D'_\eta=(G_\eta,H_\eta)=H_\eta\)\textup{)} and \(D'_s=(H_s,H_s)\) \textup{(}respectively
\(D'_s=(G_s,H_s)\)\textup{)} for \(s\in S'\); the closed subschemes \(D'\) and \(H|_{S'}\) of \(G|_{S'}\), of finite
presentation over \(S'\) and with the same generic fiber, coincide over a smaller nonempty open subset
(EGA~IV\(_3\), 8.8.2.5).}
